\documentclass[UTF8, 11pt]{ctexart}

\usepackage[margin=1in]{geometry}
\usepackage{amsmath, amssymb}
\usepackage{booktabs}
\usepackage{hyperref}
\usepackage{url}
\usepackage{enumitem}
\setlist{leftmargin=*}
\setlength{\parskip}{0.35em}
\setlength{\parindent}{0pt}

\title{\textbf{阶段 5：可选拓展}\\
启发式改进、应用层仿真与未来工作}
\author{小组成员：姓名 1，姓名 2，姓名 3}
\date{CS240：算法设计与分析，2026 年春季学期}

\begin{document}

\maketitle

\section{拓展阶段的目的}

必做部分已经实现并评估了一个经典图划分方法。本阶段描述超出最低要求的拓展
内容。目标不是声称达到最先进图划分性能，而是展示核心算法思想如何被改写、
改进，并更深入地连接到分布式优化这一现代应用场景。

本项目包含一个已经实现的拓展和若干自然的未来拓展。已经实现的拓展是轻量级
一致性平均通信仿真，它把图划分与跨工作节点 通信成本联系起来。未来拓展包括
加权通信代价、更强的数据预处理、稀疏谱方法、多层图划分，以及更真实的分布式
优化仿真。

\section{已实现拓展：一致性通信仿真}

主报告不仅使用纯图指标评估划分，还加入了一个小型分布式一致性仿真。每个顶点
持有一个标量值，相邻顶点反复交换信息。更新过程类似标准平均一致性：
\[
    x_i^{(t+1)}
    =
    x_i^{(t)}
    +
    \alpha
    \sum_{j:(i,j)\in E}
    \left(x_j^{(t)}-x_i^{(t)}\right),
\]
其中 $\alpha$ 根据最大度数选取。目标值是所有节点初始值的平均值。

划分用于统计有多少消息跨越 工作节点边界。对于端点被分配到不同分区的每条边，
仿真在每次迭代中统计两个方向的跨工作节点 消息。因此可以得到一个与割边数量
直接相关的通信指标：
\[
    \text{cross messages}
    \approx
    2T \cdot |\{(u,v)\in E:p(u)\neq p(v)\}|,
\]
其中 $T$ 是一致性迭代次数。

这一拓展为图划分目标提供了应用层解释。更低的割值意味着更少的跨工作节点 消息，
而这正是分布式系统希望降低的成本。

\section{拓展实验结果}

平均跨工作节点 消息数为：
\[
    \text{greedy}=15460,\qquad
    \text{spectral}=11360,\qquad
    \text{Kernighan--Lin}=10440.
\]
这些结果与割边权重对比一致。Kernighan--Lin 细化具有最低平均割值，因此也具有
最低平均跨工作节点 通信量。这支持了一个观点：经典图划分算法可以作为分布式
优化中的通信放置工具。

同一张图上，不同划分的一致性误差几乎相同。这是符合预期的，因为仿真保持原始
图拓扑和平均动力学不变。划分只影响消息被放置在哪里，而不改变哪些顶点相互
通信。这个结果有助于明确当前拓展测量的内容：它测量通信放置成本，而不是收敛
加速。

\section{可能拓展一：加权通信代价}

当前实验主要使用无权图或均匀权重图。更真实的分布式优化系统通常具有非均匀
通信代价。例如，某些依赖需要传输大型张量，某些 工作节点之间网络链路较慢，
某些变量之间耦合更强。

一个自然拓展是为边赋予有意义的权重 $w_{uv}$，并优化
\[
    \sum_{(u,v)\in E,\;p(u)\neq p(v)} w_{uv}.
\]
权重可以来自：
\begin{itemize}
    \item 数据依赖强度；
    \item 任务之间需要传输的字节数；
    \item 历史通信频率；
    \item 工作节点组之间的网络延迟。
\end{itemize}
这一拓展会使划分目标更接近真实分布式系统。

\section{可能拓展二：更好的预处理}

图预处理可能显著影响划分质量和运行时间。简单预处理包括删除孤立点、合并重复
或近似重复节点、归一化边权，以及在划分前按照度数或局部性对顶点排序。

对于贪心基线，更好的顶点顺序可以改善划分质量。当前实现按照加权度数排序。
更高级的顺序可以使用广度优先遍历、图退化序，或者在有空间坐标时使用空间位置。
对于谱方法，预处理还可以更谨慎地处理非连通分量，因为图存在多个连通分量时，
拉普拉斯特征向量可能不唯一或不稳定。

\section{可能拓展三：稀疏谱划分}

当前谱划分实现使用稠密特征分解，简单但不够可扩展。更强版本可以使用稀疏矩阵
和迭代特征求解器，只计算 Fiedler 向量。这会降低内存使用，并让方法适用于更大
的稀疏图。

算法改动包括：
\begin{enumerate}
    \item 将图拉普拉斯矩阵构造为稀疏矩阵；
    \item 使用迭代求解器计算最小的非零特征对；
    \item 根据 Fiedler 向量递归二分；
    \item 对很小或非连通的分量使用简单切分作为回退。
\end{enumerate}

这一拓展直接解决可扩展性问题，也会使谱基线在大图上更有竞争力。

\section{可能拓展四：多层图划分}

更进一步的拓展是实现 METIS 风格的多层算法。该算法通常包含三个阶段：
\begin{enumerate}
    \item \textbf{粗化}：反复匹配并收缩顶点，生成更小的图；
    \item \textbf{初始划分}：在最小的粗图上求一个划分；
    \item \textbf{反粗化与细化}：将划分逐层投影回更大的图，并在每一层进行
    局部细化。
\end{enumerate}

这种方法可以结合全局可扩展性和局部改进。Kernighan--Lin 或相关局部搜索方法
可以作为反粗化阶段的细化器。该拓展非常自然，因为它把本项目复现的算法与最
成功的实用图划分框架之一连接起来。

\section{可能拓展五：更强的分布式优化仿真}

当前一致性仿真统计跨工作节点 消息数，但没有让跨工作节点 通信变慢或更不可靠。
更强的拓展可以引入通信成本模型：
\begin{itemize}
    \item 跨工作节点 消息具有更高延迟；
    \item 跨工作节点 通信受到预算限制；
    \item 延迟消息会影响更新规则；
    \item 不同划分导致不同的 wall-clock 收敛时间。
\end{itemize}

例如，延迟一致性更新可以区分本地消息和远程消息。本地消息立即更新，远程消息
延迟若干轮后应用。这样，跨工作节点 边更少的划分不仅会减少通信量，也可能提升
单位时间内的收敛速度。

另一个拓展是分布式梯度下降。每个顶点可以代表一个局部目标函数 $f_i(x)$，智能体
执行局部梯度步和邻居平均。划分则决定哪些同步事件是本地的，哪些是远程的。

\section{项目层面的总结}

可选拓展进一步加强了经典算法和现代应用之间的联系。已经实现的一致性仿真表明，
更低的图割可以转化为更少的跨工作节点 消息。提出的未来拓展则可以让系统模型更
真实、算法更可扩展。

最有价值的下一步是结合加权通信模型和稀疏谱划分。这可以保留当前项目结构，
同时解决两个重要限制：边权过于均匀和稠密特征分解不够可扩展。更大的未来版本
还可以加入多层粗化，并与 METIS 风格工具直接比较。

\end{document}

